% ======================================================================
% Supplementary materials — methodological checks not central to the core claims
% Extracted from §Results 2026-04-20 to keep the main narrative focused.
% All ref labels from main text remain valid (\ref{subsec:limitations_befriends},
% \ref{subsec:conditional_process}, \ref{subsec:unified_sem}).
% ======================================================================

\appendix
\section*{Приложение S1. Дополнительные методологические проверки}
\label{supp:round3}

Содержание приложения: LPA/GMM на четырёх композитах, статус ADI как самостоятельной оси, DI vs Lie-scale, октантная карта категорий свободного поля, piecewise-линейная траектория «burning empath», network-centrality после GLasso. Каждое наблюдение поддерживает или уточняет один из основных findings §Results, но не несёт самостоятельной claim-нагрузки уровня headline-результатов.

\subsection*{S1.1 LPA / GMM на четырёх композитах}
Latent profile analysis на E, P, D, A под BIC выбирает $k=2$. Кластер 0 ($n = 1461$, «нормальный»): $E = 55.8$, $P = 42.5$, $D = 52.2$, $A = 55.1$. Кластер 1 ($n = 90$, $5.8\%$ выборки): $E = 40.2$, $P = 65.2$, $D = 55.9$, $A = 42.9$ --- именно P-полюс. Из 90 человек кластера~1, 54 попадают в квадрант «low E / high P» в 2D-проекции. То есть даже чистая mixture-модель на четырёх шкалах, без априорных меток, выделяет психопатический контур как самостоятельный естественный класс. Расширение до $k=3, 4, 5$ под BIC не даёт улучшения, на которое можно было бы опираться: $k=3$ даёт $\Delta\mathrm{BIC} = -32$, но третий кластер --- это крошечная $n = 14$ «dissociated low-A» подгруппа, структурно не новая относительно high-D подвыборки.

\subsection*{S1.2 ADI не самостоятельная ось}
Anti-Dissociation Index (ADI), рассматриваемый в части психометрической литературы как самостоятельный конструкт, в наших данных практически совпадает с $D_{\mathrm{score}}$: $r(\mathrm{ADI}, D) = +0.806$, $r(\mathrm{ADI}, \mathrm{AIS}) = +0.770$, $r(\mathrm{ADI}, \mathrm{IDS}) = +0.786$, $r(\mathrm{ADI}, P) = +0.065$ (практически ноль). ADI --- это свёртка блокировочного пояса, не отдельная латентная координата. Мы поэтому не включаем ADI как независимый предиктор в основные модели; в описательной статистике он работает как global severity index для защитного полюса, не более.

\subsection*{S1.3 DI не сводится к социальной желательности}
Discrepancy Index (разрыв между self-reported E и поведенческими маркерами) в часто задаваемом вопросе «не Lie ли это?» получает отрицательный ответ: контраст HIGH-DI vs LOW-DI подгрупп по Lie-шкале даёт $d = 0.00$ ($p = 0.99$). В совместной модели $\mathrm{DI} \sim A + P + D + \mathrm{Lie} + \mathrm{ACE}$ ($R^2 = 0.41$) Lie вносит минимальный вклад; основные драйверы --- $A$ ($+$), $P$ ($-$), $D$ ($+$). DI --- реальный структурный паттерн «аффективная эмпатия, которая не становится действенной», а не артефакт стиля ответа.

\subsection*{S1.4 Октантная карта 8/16}
Проекция категорий свободного поля на октанты (знаки $z$-оценок E, P, D, A) заполняет лишь 8 из 16 возможных октантов; 8 октантов остаются пустыми. Интересны два пустых: $++++$ (все четыре шкалы одновременно высокие) --- структурно ожидаемо пуст, так как hi-E и hi-P дают противоположные требования к архитектуре; и $+++-$ (hi-E, hi-P, hi-D, low-A) --- архетип «дружит с психопатом / эмпатический контрагент», на него приходится только $n = 3$ респондента (см. §\ref{subsec:limitations_befriends}). Остальные шесть пустых октантов --- теоретически возможные, но не представленные в выборке конфигурации; их заполнение требует таргетированного подбора подвыборок в wave~2.

\subsection*{S1.5 «Burning empath» как piecewise-линейная траектория}
Нелинейный piecewise-регрессионный анализ $\mathrm{ACE} \to E_{\mathrm{AF}}$ с оптимальным knot~$= 80$ (из 100-балльной шкалы ACE) показывает два режима: до 80 $\beta = -0.085$ ($p < 10^{-7}$), после 80 $\beta = -0.498$ ($p = 7 \times 10^{-4}$). $\Delta\mathrm{AIC} = -5.37$ против чисто линейной модели --- улучшение существует, но скромное. На бинах среднее \EAF{} монотонно падает с $57.97$ (ACE $[0, 30]$) до $51.70$ (ACE $[75, 100]$). Содержательный вывод: «сгорающий эмпат» как качественно отдельная траектория в данных \emph{не выделяется}; основной эффект CA на \EAF{} линейный, нелинейный догиб в экстремальном хвосте ACE присутствует, но не меняет общую картину.

\subsection*{S1.6 Network centrality: ACE и F --- центральные хабы}
Partial-correlation network на $\{$E, P, D, A, ACE, F, EC$\}$ после GLasso-регуляризации даёт strength centrality (сумма $|$partial $r|$): $\mathrm{ACE} = 3.43$, $F = 3.42$, $D = 3.11$, $\mathrm{EC} = 2.95$, $P = 2.60$. То есть ранний адверсивный анамнез (ACE) и компонент «F» внутри $D$-пояса (аффективная дизрегуляция) --- самые сильно связанные узлы сети. Мост между $D$-кластером и ACE проходит через $F$. Сетевая структура подтверждает каузальное предположение статьи «ACE $\to$ F $\to$ D-каскад $\to$ \EAF{}» на уровне условных зависимостей.

\subsection*{S1.7 Эмпирическая производная архетипной архитектуры}
\label{supp:empirical_archetypes}

В продуктовом слое инструмента (\texttt{cluster} field в экспортируемых данных) каждому респонденту присваивается архетипный ярлык из 12 категорий (Трикстер, Странник, Аналитик и др.). Эта декомпозиция rule-based и по конструкции декоративная; её эмпирическая воспроизводимость не была проверена в основной статье. Здесь мы тестируем вопрос напрямую: \emph{какая кластерная структура естественно поддерживается данными в 6-факторном пространстве}.

\paragraph{Метод.} На стандартизованных 6-факторных композитах (E, P, M, A, ID, CA; $N = 1465$) выполнены $k$-means кластеризации при $k = 2$--$10$. Для каждого $k$ вычислены: silhouette index, общий within-sum-of-squares, и bootstrap-стабильность через Adjusted Rand Index ($50$ итераций случайных 50\%-подвыборок).

\paragraph{Стабильность.} ARI как функция $k$:

\begin{center}
\begin{tabular}{ccc}
\toprule
$k$ & ARI mean & интерпретация \\
\midrule
2 & 0.91 & очень стабильно (P-pole vs healthy) \\
3 & 0.81 & стабильно \\
\textbf{4} & \textbf{0.72} & \textbf{стабильно (рекомендуемое разрешение)} \\
5 & 0.62 & стабильно \\
6 & 0.48 & пограничная стабильность \\
8 & 0.46 & пограничная \\
10 & 0.45 & пограничная \\
\bottomrule
\end{tabular}
\end{center}

Структура $k = 4$ выбрана как оптимальный баланс между стабильностью (ARI~$> 0.70$) и интерпретативной разрешающей способностью.

\paragraph{Четыре эмпирически устойчивых кластера.}

\begin{center}
\small
\begin{tabular}{lrrrrrrr}
\toprule
Кластер & $n$ ($\%$) & $E$ & $P$ & $M$ & $A$ & ID & CA \\
\midrule
\textbf{Свидетель} & 358 (24\%) & $+0.88$ & $-0.71$ & $-0.98$ & $+0.88$ & $-1.15$ & $-0.60$ \\
\textbf{Защищённый} & 313 (21\%) & $-1.19$ & $+1.03$ & $+1.07$ & $-0.98$ & $+1.05$ & $+0.53$ \\
\textbf{Раненый} & 453 (31\%) & $+0.09$ & $-0.52$ & $-0.11$ & $-0.41$ & $+0.35$ & $+0.65$ \\
\textbf{Оператор} & 341 (23\%) & $+0.04$ & $+0.49$ & $+0.19$ & $+0.53$ & $-0.22$ & $-0.73$ \\
\bottomrule
\end{tabular}
\end{center}

Все четыре центроида читаются как теоретически осмысленные DEIC-конфигурации: \emph{Свидетель} объединяет baseline-empath и часть integrated-survivor (§\ref{subsec:integrated}); \emph{Защищённый} соответствует secondary-psychopathy архитектуре (§\ref{subsec:primary_secondary}) с активным P-нарративом, защитной маскировкой и hi-CA + hi-ID; \emph{Раненый} --- канонический trauma-marked профиль без P-надстройки; \emph{Оператор} --- функциональный non-traumatized с конституциональным P-narrative и сохранным witnessing.

\paragraph{Сравнение с продуктовым 12-архетипным слоем.} Для каждого из 12 продуктовых архетипов вычислена доля респондентов, попадающих в доминирующий empirical cluster (k = 4). Ни один из 12 не достигает критерия consistency $\geq 70$\%:

\begin{center}
\small
\begin{tabular}{lrr}
\toprule
Продуктовый архетип & $n$ & Доля в доминирующем emp-cluster \\
\midrule
Воин & 18 & 50.0\% (мал $n$) \\
Аналитик & 289 & 33.2\% \\
Странник & 362 & 35.0\% \\
Трикстер & 377 & 28.9\% \\
Шаман & 37 & 48.6\% \\
\textit{(все остальные)} & --- & $< 50\%$ \\
\bottomrule
\end{tabular}
\end{center}

Каждый продуктовый архетип распределён по всем четырём empirical-кластерам приблизительно равномерно. Архетипное именование в текущем виде \emph{не воспроизводит} наблюдаемой структуры данных и должно интерпретироваться как декоративный narrative shorthand для product-feedback, не как эмпирически валидная типология.

\paragraph{Импликация для wave-2 и для public framing.} Для replication-ready кластерной системы инструмент должен переключиться с rule-based 12-архетипа на data-driven 4-cluster (Свидетель / Защищённый / Раненый / Оператор) с явной коммуникацией респонденту: \emph{кластер --- статистически воспроизводимый регион непрерывного 6-факторного пространства, не categorical тип личности; респондент со временем может смещаться между кластерами по мере изменения архитектуры}. Эта переформулировка снимает критику D\&D-class-style identity-typing, которой подверженно текущее labelling. Empirical cluster definitions, criteria, центроиды и caveats формализованы в \texttt{psychometrics.json:empirical\_clusters} (см.\ \texttt{empirical\_archetypes\_summary.json} в supplementary materials).
